\documentclass[11pt,a4paper]{article}
% ===================================================================
%  Vorlage: Sudoku-Arbeitsblatt mit Lösungsseite
%  1. python3 sudoku_generator.py -n 4 -s mittel --seed 1 -o sudokus.tex
%  2. pdflatex vorlage_sudoku_arbeitsblatt.tex
% ===================================================================
\usepackage[T1]{fontenc}
\usepackage{lmodern}
\usepackage[ngerman]{babel}
\usepackage[a4paper,margin=2cm]{geometry}
\usepackage{xcolor}
\usepackage{sudoku}          % NICHT zusammen mit printsudoku/sudokubundle laden

\setlength\sudokusize{7.5cm}                    % Kantenlänge des Gitters
\setlength\sudokuthickline{1.6pt}               % Blockränder
\renewcommand*\sudokuformat[1]{\LARGE\sffamily#1}

\input{sudokus}              % vom Generator erzeugt

\pagestyle{empty}
\setlength{\parindent}{0pt}

\begin{document}
{\Large\bfseries\sffamily Sudoku} \hfill Name: \rule{5cm}{0.4pt}

\smallskip
Trage die Ziffern 1 bis 9 so ein, dass jede Ziffer in jeder Zeile, jeder
Spalte und jedem $3\times3$-Block genau einmal vorkommt.

\bigskip
\begin{minipage}[t]{0.49\textwidth}\centering\textbf{A}\SudokuAufgabeA\end{minipage}\hfill
\begin{minipage}[t]{0.49\textwidth}\centering\textbf{B}\SudokuAufgabeB\end{minipage}

\bigskip
\begin{minipage}[t]{0.49\textwidth}\centering\textbf{C}\SudokuAufgabeC\end{minipage}\hfill
\begin{minipage}[t]{0.49\textwidth}\centering\textbf{D}\SudokuAufgabeD\end{minipage}

\newpage
{\Large\bfseries\sffamily Lösungen}

\bigskip
\setlength\sudokusize{6cm}
\renewcommand*\sudokuformat[1]{\large\sffamily#1}
\begin{minipage}[t]{0.49\textwidth}\centering\textbf{A}\SudokuLoesungA\end{minipage}\hfill
\begin{minipage}[t]{0.49\textwidth}\centering\textbf{B}\SudokuLoesungB\end{minipage}

\bigskip
\begin{minipage}[t]{0.49\textwidth}\centering\textbf{C}\SudokuLoesungC\end{minipage}\hfill
\begin{minipage}[t]{0.49\textwidth}\centering\textbf{D}\SudokuLoesungD\end{minipage}
\end{document}
